\chapter{Homologia singularna i charakterystyka Eulera: rachunki uzupełniające}
\label{app:homologia-singularna-euler}
\index{homologia singularna@Homologia singularna!obliczenia}
\index{charakterystyka Eulera@Charakterystyka Eulera!obliczenia}

Homologia singularna, zdefiniowana w~\ref{def:homologia-singularna},
jest dostępna dla każdej przestrzeni topologicznej, ale jej grupy
łańcuchów są zwykle nieskończenie generowane. Nawet punkt ma po jednym
singularnym sympleksie stałym w każdym wymiarze. Dlatego nie można
otrzymać charakterystyki Eulera przez naprzemienne dodawanie liczb
\emph{wszystkich} singularnych sympleksów. Gdy przestrzeń ma skończony
model komórkowy, Twierdzenie~\ref{thm:komorkowa-singularna} zastępuje
ten duży kompleks skończonym kompleksem o tych samych homologiach.
Poniżej rozpisujemy wynikające stąd rachunki, przydatne zwłaszcza
przy rozkładach powierzchni i dołączaniu uchwytów.

\section{Dwa sposoby obliczania tej samej liczby}
\label{sec:app-euler-dwa-rachunki}

Niech $X$ będzie skończonym kompleksem CW. Oznaczmy przez $c_k(X)$
liczbę jego $k$-komórek, a przez
$b_k(X)=\dim_{\Q}H_k(X;\Q)$ jego $k$-tą liczbę Betti'ego.
Homologie po prawej stronie są \emph{singularne}; porównanie
z homologią komórkową pozwala jednak policzyć je z komórek.

\begin{stwierdzenie}[Wzór Eulera dla skończonego CW]
\label{prop:app-euler-cw}
Dla skończonego kompleksu CW $X$ zachodzi
\begin{equation}\label{eq:app-euler-cw}
 \chi(X):=\sum_{k\geq0}(-1)^k c_k(X)
 =\sum_{k\geq0}(-1)^k b_k(X).
\end{equation}
W szczególności suma po komórkach nie zależy od wybranego skończonego
rozkładu CW przestrzeni $X$.
\end{stwierdzenie}
\begin{proof}
Z Uwagi~\ref{rem:homologia-wspolczynniki-przestrzeni} kompleks
komórkowy $C_k^{\mathrm{CW}}(X;\Q)$ ma dokładnie $c_k(X)$ generatorów
nad $\Q$, a jego homologia jest $H_k(X;\Q)$ na mocy
Twierdzenia~\ref{thm:komorkowa-singularna}. Kompleks jest skończony,
więc możemy zastosować Twierdzenie~\ref{thm:euler-homologia}:
naprzemienna suma wymiarów grup łańcuchów równa się naprzemiennej
sumie wymiarów ich homologii. To dokładnie~\eqref{eq:app-euler-cw}.
Jeśli wybierzemy drugi skończony rozkład CW tej samej przestrzeni,
obie sumy komórkowe są równe tej samej sumie liczb Betti'ego.
\end{proof}

Rachunek nad $\Q$ odrzuca torsję: na przykład grupa $\Z/2\Z$
wnosi zero do rangi homologii całkowitej. Można też liczyć nad dowolnym
ciałem $\mathbb F$. Wtedy ta sama liczba komórek daje
\begin{equation}\label{eq:app-euler-cialo}
 \chi(X)=\sum_k(-1)^k\dim_{\mathbb F}H_k(X;\mathbb F).
\end{equation}
Poszczególne wymiary mogą zależeć od charakterystyki ciała,
jak dla $\mathbb{RP}^2$ w Przykładzie~\ref{prz:zmiana-wspolczynnikow-rp2-torus},
lecz ich naprzemienna suma nie zależy.

\section{Pary, dołączanie komórek i sklejanie}
\label{sec:app-euler-sklejanie}

Wzór na charakterystykę względną mierzy, ile komórek dodajemy do
już zbudowanej przestrzeni. W filtracji uchwytowej odpowiednikiem
nowej komórki jest rdzeń uchwytu.

\begin{stwierdzenie}[Charakterystyka Eulera pary CW]
\label{prop:app-euler-wzgledna}
Niech $A\subseteq X$ będzie podkompleksem skończonego kompleksu CW.
Wtedy homologie względne są skończenie wymiarowe nad $\Q$ i
\begin{equation}\label{eq:app-euler-wzgledna}
 \sum_k(-1)^k\dim_{\Q}H_k(X,A;\Q)
 =\chi(X)-\chi(A).
\end{equation}
Jeśli $X$ powstaje z $A$ przez dołączenie jednej $q$-komórki,
to $\chi(X)=\chi(A)+(-1)^q$.
\end{stwierdzenie}
\begin{proof}
Długi ciąg homologii pary $(X,A)$ nad $\Q$ z
Sekcji~\ref{sec:homologia-wzgledna} jest dokładny.
Grupy $H_k(X;\Q)$ i $H_k(A;\Q)$ są skończenie wymiarowe i znikają
ponad wymiarem $X$ na mocy Twierdzenia~\ref{thm:komorkowa-singularna}.
Dokładność pokazuje, że grupy względne także są skończenie
wymiarowe i znikają w dostatecznie wysokich stopniach.
W skończonym ciągu dokładnym przestrzeni liniowych naprzemienna suma
wymiarów wynosi zero: wymiar każdego wyrazu jest sumą wymiarów
obrazu strzałki wchodzącej i wychodzącej, a te składniki znoszą się
parami. Wyrazy długiego ciągu występują w porządku
$H_k(A)\to H_k(X)\to H_k(X,A)\to H_{k-1}(A)$;
po zsumowaniu dostajemy
$0=\chi(A)-\chi(X)+\sum_k(-1)^k\dim_{\Q}H_k(X,A;\Q)$.
Zatem
$\chi(X)=\chi(A)+\sum_k(-1)^k\dim_{\Q}H_k(X,A;\Q)$,
co jest tezą. Jeśli dodajemy jedną $q$-komórkę, to liczby komórek
$X$ i $A$ różnią się tylko o~$1$ w stopniu $q$.
\end{proof}

\begin{wniosek}[Liczba uchwytów i charakterystyka Eulera]
\label{cor:app-euler-uchwyty}
Niech zwarty kobordyzm $W$ od $N_-$ ma skończony rozkład uchwytowy
z $u_k$ uchwytami indeksu $k$, a $N_-$ niech ma skończony model CW.
Wówczas
\[
 \chi(W)-\chi(N_-)=\sum_k(-1)^k u_k.
\]
\end{wniosek}
\begin{proof}
Twierdzenie~\ref{thm:kompleks-uchwytow-18} daje skończony kompleks
całkowity; po tensorowaniu z $\Q$ zgodnie z
Twierdzeniem~\ref{thm:homologia-char-zero} oblicza on
$H_*(W,N_-;\Q)$ i ma przestrzeń wymiaru $u_k$ w stopniu $k$.
Ponieważ $N_-$ ma skończony model CW,
długi ciąg pary zapewnia także skończoność sumy homologicznej
określającej $\chi(W)$ w~\eqref{eq:app-euler-cw}.
Twierdzenie~\ref{thm:euler-homologia}
utożsamia naprzemienną sumę tych wymiarów z naprzemienną sumą
wymiarów homologii względnych. Długi ciąg pary, dokładnie jak
w dowodzie Stwierdzenia~\ref{prop:app-euler-wzgledna}, utożsamia
tę ostatnią sumę z $\chi(W)-\chi(N_-)$.
\end{proof}

\begin{stwierdzenie}[Addytywność przy sklejaniu podkompleksów]
\label{prop:app-euler-addytywnosc}
Jeśli $X=A\cup B$, gdzie $A$ i $B$ są podkompleksami skończonego
kompleksu CW $X$, to
\begin{equation}\label{eq:app-euler-sklejanie}
 \chi(X)=\chi(A)+\chi(B)-\chi(A\cap B).
\end{equation}
\end{stwierdzenie}
\begin{proof}
Każda otwarta komórka $X$ należy do $A$, do $B$ albo do obu.
W każdym wymiarze zwykłe liczenie zbiorów komórek daje
$c_k(X)=c_k(A)+c_k(B)-c_k(A\cap B)$.
Mnożymy tę równość przez $(-1)^k$ i sumujemy po skończonej liczbie
wymiarów. Wzór~\eqref{eq:app-euler-cw} zapewnia, że wynik jest też
rachunkiem z homologii singularnej.
\end{proof}

\begin{przyklad}[Sfera z dwóch dysków]
Rozłóżmy $S^2$ na dwa domknięte dyski ze wspólnym brzegiem $S^1$
i wybierzmy zgodne, skończone struktury CW, dla których dyski oraz
okrąg są podkompleksami. Dysk ma jedną $0$-komórkę, jedną
$1$-komórkę na brzegu i jedną $2$-komórkę, więc $\chi(D^2)=1$;
okrąg ma po jednej komórce w wymiarach $0$ i $1$, więc
$\chi(S^1)=0$. Wzór~\eqref{eq:app-euler-sklejanie} daje
$\chi(S^2)=1+1-0=2$. Zgadza się to z obliczeniem
$H_0(S^2;\Q)=H_2(S^2;\Q)=\Q$ i $H_1(S^2;\Q)=0$
z Twierdzenia~\ref{thm:homologia-sfer}.
\end{przyklad}

\section{Wielokąty powierzchni: brzegi i cecha Eulera}
\label{sec:app-euler-powierzchnie}

Rozważmy najpierw powierzchnię $\Sigma_g$ otrzymaną ze standardowego
wielokąta o słowie brzegowym
\[
 a_1b_1a_1^{-1}b_1^{-1}\cdots
 a_gb_ga_g^{-1}b_g^{-1},\qquad g\geq1.
\]
Po sklejeniu wszystkie wierzchołki dają jeden wierzchołek,
pary boków dają $2g$ krawędzi, a wnętrze wielokąta jedną
$2$-komórkę. Obliczmy brzegi, zamiast zatrzymywać się na samym
liczeniu komórek. Każda $1$-komórka jest pętlą, więc
$\partial_1=0$. Współczynnik $\partial_2$ przy danej
$1$-komórce jest sumą jej przebiegów w słowie brzegowym ze znakami
orientacji, zgodnie z konstrukcją w~\ref{sec:homologia-komorkowa}.
Każda litera pojawia się tu raz dodatnio i raz ujemnie, zatem
$\partial_2=0$. Kompleks i jego homologie mają więc postać
\[
 0\longrightarrow\Z\xrightarrow{\;0\;}\Z^{2g}
 \xrightarrow{\;0\;}\Z\longrightarrow0,
 \qquad
 H_2(\Sigma_g;\Z)=\Z,\quad
 H_1(\Sigma_g;\Z)=\Z^{2g},\quad
 H_0(\Sigma_g;\Z)=\Z.
\]
Stąd $\chi(\Sigma_g)=1-2g+1=2-2g$.
Dla $g=0$ analogiczny model sfery ma jedną $0$-komórkę
i jedną $2$-komórkę; ten sam wzór daje $\chi(S^2)=2$.

Model nieorientowalny $N_h$ otrzymujemy z wielokąta o słowie
$c_1c_1\cdots c_hc_h$, gdzie $h\geq1$.
Znów jest jeden wierzchołek i jedna $2$-komórka, lecz tym razem
jest $h$ krawędzi, a każdą z nich brzeg obiega dwa razy w tym samym
kierunku. Dla standardowej bazy $e_1,\ldots,e_h$ mamy zatem
\[
 0\longrightarrow\Z\xrightarrow{\;\partial_2\;}\Z^h
 \xrightarrow{\;0\;}\Z\longrightarrow0,
 \qquad \partial_2(1)=2(e_1+\cdots+e_h).
\]
Wektor $v=e_1+\cdots+e_h$ można uzupełnić do bazy
$(v,e_2,\ldots,e_h)$ grupy $\Z^h$, ponieważ
$e_1=v-e_2-\cdots-e_h$.
Dlatego $\partial_2$ jest injekcją i
\[
 H_2(N_h;\Z)=0,\qquad
 H_1(N_h;\Z)\cong\Z^{h-1}\oplus\Z/2\Z,\qquad
 H_0(N_h;\Z)=\Z.
\]
Cecha Eulera wynosi $\chi(N_h)=1-h+1=2-h$;
po przejściu do współczynników $\Q$ tę samą wartość daje
$1-(h-1)$. Przypadek $h=1$ odtwarza rachunek dla
$\mathbb{RP}^2$ z Przykładu~\ref{prz:homologia-rp2}.
Te obliczenia dotyczą wskazanych modeli wielokątowych; pełne
twierdzenie o klasyfikacji wszystkich zwartych powierzchni
jest osobnym wynikiem używanym w Rozdziale~\ref{ch:poincare-2-high}.
